Pour mesurer la largeur d'une fente, on l'éclaire avec un laser de longueur d'onde \qty{\v{lam}}{nm}. Sur un écran placé à \qty{\v{L}}{m}, le maximum central de diffraction mesure \qty{\v{W}}{mm} de large.

\begin{enumerate}[label=\alph*)]
    \item Déterminez la largeur de la fente.
    \item À quel angle se trouve le troisième minimum de diffraction ?
    \item Combien de minimums de diffraction peut-on observer de chaque côté du maximum central (en supposant un écran infiniment grand) ?
\end{enumerate}
